% Session 10 -- the neural network as a parametric function, its training loop, and the four
% places the very same template is reused later in the session.
% Placement: sec. "Neural networks as function approximators", after equation (10.1) and the
% cat-versus-dog paragraph.
\begin{aqmfigure}{fig:s10:template}{The construction that the rest of this session uses again and again.
\textbf{Left.} A deep network is a parametric function $f_\theta$. The architecture, which is the
number of layers and their width, fixes which functions it \emph{can} represent, and the weights
and biases $\theta$ fix which one it \emph{does} represent at a given moment. \textbf{Right.}
Training is the loop of (10.1). We evaluate the output, compare it with the label through a loss,
obtain $\partial\mathcal{L}/\partial\theta$ exactly with backpropagation, take a gradient step and
repeat. \textbf{Bottom.} Every later construction in this session is the same loop with a different
choice of input, output and loss, and nothing else changes. Each row should be read together with
the section that it names, since the physics is only in the three columns and never in the
training method.}

\begin{tikzpicture}[
  unit/.style={aqm, circle, draw=black, fill=white, minimum size=5.5pt, inner sep=0pt},
  unitin/.style={unit, fill=aqmfill, draw=aqmaccent},
  unitout/.style={unit, fill=aqmfill, draw=aqmaccent, minimum size=7pt},
  wire/.style={draw=aqmmuted, line width=0.22pt},
  tbox/.style={aqmbox, text width=5.0cm, font=\footnotesize, anchor=north west},
  hdr/.style={font=\footnotesize\bfseries, anchor=north west, inner sep=1pt},
  cell/.style={font=\footnotesize, anchor=north west, inner sep=1pt, align=left},
]

% ===================== LEFT: the network itself ==============================
\draw[aqmcell] (0.30,0.55) rectangle (7.95,3.95);
\node[aqmlabsmall, anchor=south west, text=aqmmuted] at (0.42,0.60) {$f_\theta$};

% input thumbnail
\begin{scope}[shift={(0.70,1.90)}]
  \foreach \i in {0,1,2,3} \foreach \j in {0,1,2,3}
     \draw[aqm, draw=aqmmuted, line width=0.3pt, fill=aqmfill]
        (\i*0.16,\j*0.16) rectangle (\i*0.16+0.16,\j*0.16+0.16);
\end{scope}
\node[aqmlabsmall, anchor=north, align=center, text width=1.2cm] at (1.02,1.68) {$N\times N$\\[-2pt] image};

% layers
\foreach \j in {0,...,4} \node[unitin] (I\j) at (2.05,0.42*\j+1.32) {};
\foreach \j in {0,...,5} \node[unit]   (A\j) at (3.35,0.42*\j+1.11) {};
\foreach \j in {0,...,5} \node[unit]   (B\j) at (4.65,0.42*\j+1.11) {};
\node[unitout] (O0) at (5.95,2.46) {};
\node[unitout] (O1) at (5.95,1.86) {};

\begin{scope}[on background layer]
  \foreach \a in {0,...,4} \foreach \b in {0,...,5} \draw[wire] (I\a) -- (A\b);
  \foreach \a in {0,...,5} \foreach \b in {0,...,5} \draw[wire] (A\a) -- (B\b);
  \foreach \a in {0,...,5} \foreach \b in {0,1}     \draw[wire] (B\a) -- (O\b);
\end{scope}

\draw[aqmarrow, line width=0.6pt] (1.45,2.22) -- (1.86,2.22);

\node[aqmlabsmall, anchor=north] at (2.10,1.00) {input $x$};
\node[aqmlabsmall, anchor=north] at (3.90,1.00) {hidden layers};
\node[aqmlabsmall, anchor=north] at (6.15,1.00) {output $f_\theta(x)$};
\node[aqmlabsmall, anchor=west] at (6.20,2.46) {``cat'' $(1,0)$};
\node[aqmlabsmall, anchor=west] at (6.20,1.86) {``dog'' $(0,1)$};

\node[aqmlabsmall, anchor=north, align=center, text=aqmmuted] at (4.10,3.90)
   {one weight per wire, one bias per node: together $\theta$};

% ===================== RIGHT: the training loop ==============================
\node[aqmlabsmall, anchor=north west] (tset) at (9.55,4.85)
   {labelled training set $\{(x_n,y_n)\}$};
\node[tbox] (loss) at (9.55,4.05)
   {$\mathcal{L}(\theta)=\sum_n\big\|f_\theta(x_n)-y_n\big\|^2$\ \ (10.1)\\[1pt]
    how far the output sits from the label};
\node[tbox] (grad) at (9.55,2.55)
   {\textbf{backpropagation}: $\partial\mathcal{L}/\partial\theta_k$ for every parameter at
    once, exactly, by the chain rule};
\node[tbox, draw=aqmaccent, fill=aqmfill] (upd) at (9.55,1.15)
   {$\theta \longrightarrow \theta - \eta\,\nabla_\theta\mathcal{L}(\theta)$\\[1pt]
    one gradient-descent step, learning rate $\eta$};

\draw[aqmarrow] (tset.south) -- (loss.north -| tset.south);
\draw[aqmarrow] (7.95,2.16) -- (8.75,2.16) -- (8.75,3.55) -- (loss.west);
\node[aqmlabsmall, anchor=north west] at (8.02,2.10) {$f_\theta(x_n)$};
\draw[aqmarrow] (loss.south -| 10.40,0) -- (grad.north -| 10.40,0);
\draw[aqmarrow] (grad.south -| 10.40,0) -- (upd.north -| 10.40,0);

\draw[aqmarrowaccent] (upd.south -| 10.40,0) -- (10.40,-0.62) -- (3.90,-0.62) -- (3.90,0.55);
\node[aqmlabsmall, anchor=south, text=aqmaccent, fill=white, inner sep=2.5pt] at (7.15,-0.60)
   {every weight and bias is nudged, and the whole loop runs again};

% ===================== FOOTER: the same loop, four times =====================
\draw[aqm, draw=black, line width=0.6pt] (0,-1.35) -- (15.0,-1.35);
\node[anchor=south west, font=\small\bfseries] at (0,-1.27)
   {The same loop, four more times in this session --- only the three columns change};

\draw[aqm, draw=aqmmuted, line width=0.4pt] (0,-2.05) -- (15.0,-2.05);
\node[hdr] at (0.00,-1.45) {construction};
\node[hdr] at (4.60,-1.45) {input};
\node[hdr] at (7.85,-1.45) {output};
\node[hdr] at (11.20,-1.45) {loss, minimized over $\theta$};

\foreach \y/\a/\b/\c/\d in {%
  -2.15/{NNQS ground state \S10.3}/{a spin configuration $s$}/{the coefficient $C_\theta(s)$}/{the variational energy $E[\theta]$, (10.3)},
  -3.00/{state tomography \S10.3}/{a spin configuration $s$}/{modulus and phase of $C_\theta(s)$}/{$-\sum_s n(s)\ln p_\theta(s)$},
  -3.85/{phase classification \S10.6}/{a spin snapshot}/{a one-hot phase label}/{(10.1), entirely unchanged},
  -4.70/{Hamiltonian learning \S10.7}/{an observable $\mathcal{O}$}/{the parameters $\lambda$}/{$\sum_n\|\mathcal{N}_\theta(\mathcal{O}_n)-\lambda_n\|^2$, (10.12)}%
} {
  \node[cell, text width=4.35cm] at (0.00,\y) {\a};
  \node[cell, text width=3.10cm] at (4.60,\y) {\b};
  \node[cell, text width=3.20cm] at (7.85,\y) {\c};
  \node[cell, text width=3.85cm] at (11.20,\y) {\d};
}
\draw[aqm, draw=black, line width=0.6pt] (0,-5.60) -- (15.0,-5.60);

\end{tikzpicture}

\end{aqmfigure}
